%------------------------------------------------------------------------------
% \PART 4
%------------------------------------------------------------------------------
\chapter[Разработка структуры IP-блока нейронной сети для распознавания рукописных цифр]
{РАЗРАБОТКА СТРУКТУРЫ IP-БЛОКА НЕЙРОННОЙ СЕТИ ДЛЯ РАСПОЗНАВАНИЯ РУКОПИСНЫХ ЦИФР}

%------------------------------------------------------------------------------
% \PART 4.1 Text
%------------------------------------------------------------------------------
\section{Описание структурной схемы}\par
\hspace*{12.5 mm}Принцип работы данной нейронной сети основан на комбинации 
обучаемого двумерного разделимого преобразования и полносвязной нейронной сети. 
Такая архитектура позволяет эффективно обрабатывать изображения, снижая 
вычислительные затраты при аппаратной реализации.

Для окончательной классификации используется полносвязный слой с функцией 
активации {softmax}. Данный слой преобразует выходные значения в вероятностное 
распределение по классам, позволяя определить, к какой цифре относится входное 
изображение.

Вычислительное ядро разработанного устройства состоит из 28 MAC 
(Multiply-ACcumulate) ядер, распределенных по специализированным блокам. Эти 
блоки организованы таким образом, чтобы обеспечить эффективное выполнение 
операций умножения вектора изображения на матрицу весов.

Десять MAC-ядер из общего массива используются на финальном тапе вычислений и 
непосредственно участвуют в формировании входных данных для функции активации 
softmax. Эти ядра получают входные данные из трех независимых блоков памяти, 
что позволяет оптимизировать процесс выборки весов и повысить скорость 
вычислений. Остальные 18 MAC-ядер распределены по блокам, где каждая 
вычислительная единица получает данные из двух запоминающих устройств.

IP-блок устройства предназначен для интеграции в систему на кристалле (СнК) и 
взаимодействует с процессорной системой через AXI-uP интерфейсы. СнК – это 
интегральная схема, которая объединяет в одном чипе все основные компоненты
вычислительной системы: процессор, ОЗУ и ПЗУ , графический процессор, 
контроллеры периферии. Общая структура представлена на 
рисунке~\ref{fig:str_all}.

\insertfigure{str_all}{pic/str.png}{Пример использования IP-блока нейронной сети}{16cm}

Последовательность обработки данных организуется следующим образом:

  1 Исходное изображение размером \(28\times28\) пикселей поступает на вход 
нейронной сети через переходник интерфейсов AXI-uP. Данные интерфейсы 
используется для взаимодействия с процессорной системой и обеспечивают передачу 
изображения в IP-блок. Пиксели изображения принимаются последовательно и 
записываются во внутреннюю память IP-блока, предназначенную для хранения 
изображения на протяжении всего цикла обработки.

  2 Каждый отдельный пиксель изображения поочередно передаётся на вход массива 
из 28 умножающих-накопительных ядер (MAC), которые одновременно выполняют 
операции умножения входного значения пикселя на соответствующий весовой 
коэффициент. Весовые коэффициенты извлекаются из специализированных блоков 
памяти весов, которые хранят параметры, полученные в результате 
предварительного обучения модели.

  3 За согласованность и последовательность операций отвечает управляющий 
модуль. Он обеспечивает адресацию ячеек памяти изображения и весов, организует 
подачу данных в вычислительные блоки и отслеживает фазы обработки.

  4 После завершения расчетов в LST-1 слое, результат, новое 
представление входного изображения, передаётся во второй уровень вычислений. 
Здесь используются 10 MAC-ядер, каждое из которых обрабатывает соответствующий 
вектор признаков, поступающий от первого слоя. В результате получается выходной 
вектор размерности 10, где каждый элемент отражает степень принадлежности 
изображения к соответствующему классу (от 0 до 9).

  5 Полученный выходной вектор подаётся в блок активации softmax. 
На выходе формируется распределение 
вероятностей принадлежности изображения к каждому из 10 классов.

  6 Индекс элемента вектора softmax, имеющий наибольшее значение, определяется и передаётся через интерфейс AXI4-lite 
обратно в процессорную систему. Интерфейс uP (user processor 
interface) позволяет интегрировать модель в более широкую систему на кристалле, 
облегчая её взаимодействие с внешним программным обеспечением и обеспечивая 
гибкость конфигурации \cite{chu_fpga}.

  Временная диаграмма работы uP-интерфейса представлена на 
рисунке~\ref{fig:uP}.

\insertfigure{uP}{pic/up.png}{Временная диаграмма транзакций чтения/записи интерфейса uP}{12.5cm}

  AXI4-Lite – это упрощённый вариант AXI4-интерфейса, предназначенный для 
управления и обмена небольшими порциями данных. В отличие от AXI4, он 
поддерживает только однослотовые транзакции без разделения на несколько 
пакетов\cite{vivado_axi}.

  Временная диаграмма работы AXI4-lite-интерфейса представлена на 
рисунке~\ref{fig:axi4}.
    
\insertfigure{axi4}{pic/axi.png}{Временная диаграмма транзакций чтения/записи интерфейса AXI4-lite}{13cm}

%------------------------------------------------------------------------------
% \PART 4.2 Text
%------------------------------------------------------------------------------
\section{Программная реализация эталонной модели нейронной сети на языке Python}
\hspace*{12.5 mm}Программная реализация нейронной сети выполняется с 
использованием языка Python.  
В качестве эталонной модели реализуется нейронная сеть на основе библиотеки 
{fixpoint}, позволяющая провести сравнение точности вычислений при 
использовании чисел с фиксированной запятой и чисел с плавающей запятой.

Эталонная модель представляет собой нейросеть с двумя скрытыми слоями и 
активационной функцией {tanh}.  

Для аппаратной реализации на FPGA и эталонной модели числа представляются в 
формате с фиксированной запятой (Fixed-Point). В отличие от чисел с плавающей 
запятой (Floating-Point), данный формат обеспечивает более эффективное 
использование аппаратных ресурсов, таких как блоки DSP и BRAM.\@

В данной работе используется фиксированное представление с разрядностью $Q6.7$,
где:

  – 6 бит отведены под целую часть числа.

  – 7 бит используются для дробной части.

Таким образом, числа в этом формате могут представлять значения в диапазоне 
$[-32, 31.996]$ с шагом $2^{-7} = 0.0078125$.

При переводе значений из плавающей запятой в фиксированную используется метод 
округления к ближайшему минимальному числу. 
Этот метод позволяет минимизировать ошибки округления и избежать 
систематического смещения, которое может возникать при округлении к ближайшему 
числу.

Входные изображения размером \(28 \times 28\) проходят последовательную 
обработку:

  – Первый слой выполняет линейное преобразование входных данных с весами и 
смещением, после чего применяется функция активации тангенс.
    
  – Второй слой аналогично выполняет линейное преобразование, используя 
параметры весов и смещений, после чего применяется функция активации тангенс.
    
  – Выходной слой выполняет линейное преобразование входных данных с весами и 
смещением, а затем классифицирует изображение, вычисляя вероятности классов с 
помощью softmax.

Сравнение выходных значений между моделью с фиксированной точностью и моделью с
плавающей точкой представлено на рисунке~\ref{fig:dif}.  

\insertfigure{dif}{pic/difpng.png}{Сравнение результатов вычислений}{13cm}

Графический анализ разности значений в обработке одного и того же изображения 
позволяет оценить влияние округлений и возможные потери точности. 

Разработанная 
эталонная модель на Python позволяет проверить корректность вычислений перед 
аппаратной реализацией, а также провести тестирование точности чисел с 
фиксированной запятой.

%------------------------------------------------------------------------------
% \PART 4.3 Text
%------------------------------------------------------------------------------
\section{Разработка операционной части нейронной сети}
\hspace*{12.5 mm}Нейронная сеть есть объединение управляющего и операционного
автоматов, что показано на рисунке~\ref{fig:operat}. Такое разбиение 
соответствует классическому подходу к построению цифровых вычислительных 
устройств, при котором логика работы системы делится на управление и выполнение.

\insertfigure{operat}{pic/operational.png}{Представление нейронной сети в виде управляющего и операционного автоматов}{16cm}

IP-блок нейронной сети должен быть спроектирован как независимый компонент 
системы, обеспечивающий модульность и возможность интеграции в различные 
вычислительные среды. Для этого необходимо, чтобы блок обладал четко 
определенным и простым интерфейсом, что позволит эффективно взаимодействовать 
с другими компонентами системы.

Интерфейс IP-блока должен обеспечивать передачу входных данных, управление 
процессом вычислений и выдачу результатов.

На рисунке~\ref{fig:interface} представлен требуемый вид интерфейса 
IP-блока, демонстрирующий основные входные и выходные сигналы.

Входные сигналы включают:

– clk – тактовый сигнал, синхронизирующий работу IP-блока;

– rst\_n – активный по низкому уровню сигнал сброса, приводящий систему в исходное 
состояние;

– start\_i – управляющий сигнал запуска обработки входных данных;

– addr\_i – адресный сигнал, указывающий на текущую ячейку памяти изображения;

– data\_i – входные данные, представляющие собой значение пикселя изображения, 
поступающего в IP-блок.

Выходные сигналы включают:

– num\_o – распознанные значение цифры (от 0 до 9);

– rdy\_o – сигнал готовности результата.

\insertfigure{interface}{pic/interface.png}{Интерфейс IP-блока нейронной сети}{7cm}

%------------------------------------------------------------------------------
% \PART 4.4 Text
%------------------------------------------------------------------------------
\section{Проектирование функциональной схемы нейронной сети}
\hspace*{12.5 mm}При разработке IP-блока нейронной сети необходимо определить 
его функциональную структуру, обеспечивающую корректное выполнение вычислений, 
взаимодействие с внешними компонентами системы и эффективное управление 
процессом обработки данных.

% Также на функциональной схеме показаны микрооперации и логические условия. 
Данная нейронная сеть состоит из нескольких основных блоков. Существует два 
типа блоков обработки пикселя, которые отличаются количеством блоков памяти,
которые должны обрабатываться мультиплексором и поступать в качестве входных 
данных на умножитель.
После вычисления полносвязного слоя для конкретной строки или столбца данные 
поступают в регистр обработки, состоящий из 28 чисел разрядностью 13 бит. 
Затем данные из этого регистра последовательно поступают в регистр данных 
тангенса, из которого они отправляются в блок аппроксимации тангенса 
гиперболического.
Также используются два пяти разрядных и один десяти разрядный счетчики, которые
используются для генерации адреса в память весов и в память изображения.
Блок функции активации softmax реализован на принципе сравнения всех десяти 
входных классов и выбора наибольшего значения, что будет соответствует наиболее
вероятному значению. В результате на выход поступает индекс максимального 
элемента.
Для управления работой и генерирования микроопераций на основе логических 
условий используется микропрограммный автомат.

Электрическая функциональная схема IP-блока нейронной сети прямого 
распространения приведена на рисунке~\ref{fig:LST-1-IP-block}.

\insertfigure{LST-1-IP-block}{pic/LST-1-IP-block.png}{функциональная схема IP-блока нейронной сети}{16cm}

Электрическая функциональная схема блоков обработки пикселей приведена 
на рисунке~\ref{fig:PEs}.

\insertfigure{PEs}{pic/PEs.png}{Устройство процессорных элементов}{12cm}

%------------------------------------------------------------------------------
% \PART 4.6 Text
%------------------------------------------------------------------------------
\section{Проектирование управляющей части}
\hspace*{12.5 mm}Проектирование управляющей части вычислительного устройства 
выполняется с использованием двух основных структурных компонентов:

1 Память состояний автомата, в которой хранится заданная последовательность 
управляющих слов. 

2 Комбинационная схема управления, которая отвечает за генерацию адреса 
следующего состояния автомата. Адрес определяется на основе текущего состояния 
автомата и внешних условий. 

На рисунке~\ref{fig:cnn-pic} представлена обобщённая структурная схема 
такого управляющего автомата. 

\insertfigure{cnn-pic}{pic/cnt.png}{Представление управляющей части нейронной сети в виде композиции памяти и комбинационной схемы}{10cm}

Память автомата состоит из заранее выбранных элементов памяти, обеспечивающих 
хранение информации о текущем состоянии устройства. В качестве элементов памяти 
используются триггеры, которые позволяют надежно фиксировать и изменять 
состояние системы.  

Количество элементов памяти, необходимых для хранения состояний автомата, 
определяется по следующей формуле:
\begin{equation}
    N = \log_{2}{M},
\end{equation}

\noindentгде $M$ – число состояний микропрограммного автомата.  

Каждое состояние автомата описывается представлением из шести бит в коде Грея, 
что позволяет повысить надежность работы системы, уменьшить количество 
переключений триггеров.

Дополнительно, в разработке автомата управления применяется стратегия, при 
которой старший бит кодового слова указывает на нахождение автомата в состоянии
инициализации. Это состояние используется для сброса 
счетчиков и позволяет снизить вероятность ошибок, связанных с гличами, так как для перехода
только в данное состояние необохимо установка старшего бита в 0.

Граф переходов управляющего автомата позволяет представить 
алгоритм в соответствии с условными обозначениями, что представлено 
на рисунке~\ref{fig:gsa}.

\insertfigure{gsa}{pic/gsa.png}{Граф переходов}{14cm}

Состояния граф-схемы алгоритма объеденины в блоки. Опишем каждый из них и
покажема закодированное представление:

  1 {INIT} ('b000000) – инициализация, установка начальных параметров.

  2 {LOAD} ('b100000) – загрузка изображения.

  3 {LD\_MEM} ('b100010) – подготовка данных.

  4 {INIT\_CALC\_SYNC} ('b100110) – синхронизация смещений.

  5 {INIT\_SYNC} ('b101110) – синхронизация вычислений.

  6 {INIT\_CALC} ('b101100) – инициализация смещения.

  7 {STR\_CALC} ('b111100) – обработка строки/столбца.

  8 {STR\_SYNC} ('b111110) – синхронизация обработка последнего пикселя.

  9 {STR\_RDY} ('b110110) – готовность результата обработки.

  10 {STR\_RDY\_SYNC} ('b110010) – синхронизация флага готовности.

  11 {TANG\_STRING\_CALC} ('b110000) – вычисление тангенса.

  12 {UPD\_RCO\_TYPE} ('b111000) – смена слоя.

  13 {UPD\_RCO\_SYNC} ('b111010) – синхронизация смены слоя.

  14 {INIT\_OUT\_SYNC} ('b101010) – синхронизация выходных смещений.

  15 {INIT\_OUT} ('b101011) – инициализация выходных смещений.

  16 {OUT\_CALC} ('b100011) – обработка выходного полносвязного слоя.

  17 {OUT\_RDY} ('b100111) – готовность детектирования.

  